\documentclass[11pt]{article}
\usepackage[a4paper,top=19mm,bottom=19mm,left=24mm,right=24mm]{geometry}
\usepackage[T1]{fontenc}
\usepackage{newtxtext}
\usepackage{microtype}
\usepackage[hidelinks]{hyperref}
\pagestyle{empty}
\setlength{\parindent}{0pt}
\setlength{\parskip}{5pt}

\begin{document}
\raggedright
\hyphenpenalty=10000

25 July 2026

Editorial Office\\
\textit{Expert Systems with Applications}

\textbf{Re: Research article submission, ``ContractRAG: A
risk-certified expert system for deploying retrieval-augmented
generation services''}

Dear Editor,

On behalf of all authors, we submit this manuscript for consideration
as a Research Article in \textit{Expert Systems with Applications}.

Organizations deploying retrieval-augmented generation (RAG) must
choose among many retrieval, reranking, and generation configurations.
Current systems usually recommend a configuration from average
development-set performance. They do not tell an operator how often the
deployed service may return a wrong, unsupported, stale, or late
answer. Our manuscript presents ContractRAG, a risk-aware expert system
that converts such operational requirements into an auditable
deployment decision.

A domain expert first defines an evidence contract containing measurable
losses, admissible violation rates, and a confidence budget. ContractRAG
then compiles heterogeneous knowledge-access plans, learns progressive
policies that acquire stronger evidence only when needed, and uses a
distribution-free test to recommend the least-cost policy in a
certified prefix. When no candidate can meet the contract, the system
reports infeasibility instead of silently weakening the requirement.
After deployment, an anytime-valid monitor detects drift and triggers a
fresh decision cycle. The result is a complete design--test--deploy--monitor
workflow rather than a stand-alone routing component.

We evaluate this workflow on HybridQA, CRAG, ASQA, and QAMPARI using
four large language model families and two serving backends. Across
1,000 calibration draws, recommended policies exceed their contracts
in at most 0.8\% of draws under a 10\% error budget, compared with up to
64\% for empirical and Bayesian tuning. They reduce model cost by up to
20$\times$, or 10--18$\times$ after audit cost. The evaluation also
covers strict and joint contracts, human deferral, infeasible
requirements, audit sampling, and post-deployment drift.

The work fits the journal's focus on the design, development, testing,
and management of expert and intelligent systems. It supplies an
operator-facing decision process for information retrieval and risk
assessment, an implemented architecture, formal guarantees, and
practical evidence across several public tasks.

This manuscript is original, has not been published previously, and is
not under consideration elsewhere. All authors have approved the
manuscript and its submission. The authors declare no competing
financial interests or personal relationships that could have
influenced the work.

Thank you for considering our submission.

Sincerely,

\textbf{Weiwei Fu and Jinjiang Cui}\\
Co-corresponding authors, on behalf of all authors\\
Suzhou Institute of Biomedical Engineering and Technology\\
Chinese Academy of Sciences\\
No.\ 88 Keling Road, Suzhou New District, Suzhou 215163, China\\
\href{mailto:fuww@sibet.ac.cn}{fuww@sibet.ac.cn};
\href{mailto:cuijj@sibet.ac.cn}{cuijj@sibet.ac.cn}

\end{document}
